\documentclass[12pt]{article}
\usepackage[utf8]{inputenc}
\usepackage[T1]{fontenc}
\usepackage{hyperref}
\title{Pulabu Language Resources for K-12 Curriculum Development}
\author{Compiled by Claude Sonnet 5.5 (claude-sonnet-5-5) for Andr\'as Kornai}
\date{}
\begin{document}
\maketitle

\section{Introduction}
Pulabu (ISO 639-3 \texttt{pup}, Glottocode pula1267) is a Rai Coast (Kabenau) language of Madang Province, Papua New Guinea, in the Trans-New Guinea family. Wikipedia reports about 120 speakers from a 1981 figure. Search found almost nothing beyond encyclopedic and catalogue entries.

\section{Primary Educational Resources}
No primers, textbooks, or school programmes were found.

\section{Linguistic and Reference Materials}
\begin{itemize}
\item Pulabu language (Wikipedia): \url{https://en.wikipedia.org/wiki/Pulabu_language}
\item Glottolog language page (map view): \url{https://glottolog.org/resource/languoid/id/pula1267.bigmap.html}
\item Endangered Languages Project entry: \url{https://endangeredlanguages.com/node/30428}
\end{itemize}

\section{Religious and Community Resources}
Joshua Project lists the Pulabu people; the page is generic and not linguistic: \url{https://joshuaproject.net/people_groups/14470}

\section{Digital and Audio Resources}
None found.

\section{Curriculum Development Context}
No documentation suitable for curriculum use was found. Any programme would require new documentation work with the community, with Tok Pisin and English as the likely school languages.

\end{document}